\subsection{نتایج تجربی}
نتایج حاصل از ارزیابی تجربی بنچ‌مارک یکپارچه برای ماژول \lr{slicegpt\_pca} بر روی مدل‌های زبانی هدف در جدول \ref{tab:benchmark_slicegpt_pca} گزارش شده است:

\begin{table}[H]
\centering
\small
\caption{نتایج بنچ‌مارک تجربی ماژول فشرده‌ساز \lr{slicegpt\_pca} روی معماری‌های گوناگون}
\label{tab:benchmark_slicegpt_pca}
\begin{tabular}{lccccc}
\toprule
\textbf{مدل پایه \lr{(Base Model)}} & \textbf{دقت \lr{(Accuracy)}} & \textbf{شاخص \lr{F1}} & \textbf{حافظه \lr{(Peak VRAM)}} & \textbf{زمان ارزیابی \lr{(Time)}} & \textbf{وضعیت} \\
\midrule
\lr{DistilBERT} & \lr{90.48\%} & \lr{0.9077} & \lr{312.2 MB} & \lr{25.9 s} & \lr{PASS} \\
\lr{RoBERTa} & \lr{67.66\%} & \lr{0.7496} & \lr{301.0 MB} & \lr{30.9 s} & \lr{PASS} \\
\lr{GPT-2} & \lr{62.04\%} & \lr{0.6783} & \lr{637.3 MB} & \lr{37.1 s} & \lr{PASS} \\
\lr{Qwen} & -- & -- & -- & -- & \lr{NOT\_RUN} \\
\bottomrule
\end{tabular}
\end{table}

بر اساس نتایج تجربی، مدل \lr{DistilBERT} دقت پایه خود (\lr{90.48\%}) را به طور کامل و بدون کوچکترین افت حفظ کرده است (\lr{Delta Acc: +0.00\%}). در مدل خودبازگشتی \lr{GPT-2} دقت با رشد مواجه شده و به \lr{62.04\%} رسیده است (\lr{Delta Acc: +1.26\%}). در مدل \lr{RoBERTa} به دلیل ساختار پیچیده فضاهای نهفته لایه‌های میانی، افت دقت \lr{-7.45\%} مشاهده شده اما حدنصاب پذیرش با وضعیت \lr{PASS} احراز گردیده است. مدل \lr{Qwen} به دلیل الزامات محاسباتی لایه‌های چندگروهی (\lr{GQA}) در چرخه بنچ‌مارک به صورت \lr{NOT\_RUN} ثبت شده است.
